\section{Introduction}
Activation steering has emerged as a compelling alternative to fine-tuning for shaping the behavior of large language models (LLMs) at inference time \citep{rimsky2024steering, li2023inference, wang2025truthflow}. Its empirical success is often motivated by the linear representation hypothesis \citep{park2023linear, alain2016understanding}, where many high-level properties encoded by an LLM, such as sentiment, truthfulness, or harmfulness, appear to correspond to approximately linear directions in activation space. Under this view, changing a model's behavior can be achieved by moving its internal activations along directions associated with the target property at inference time. A standard approach is activation addition, which modifies the original activations by adding either a fixed steering vector or one that depends on the current activation.

This perspective has inspired a growing body of work on inference-time control through activation interventions. Early methods typically construct a single steering direction from contrastive examples and then apply that direction uniformly across queries \citep{rimsky2024steering}, while more recent approaches refine the dynamics of the intervention by integrating a continuous steering field over multiple steps \citep{zhaoodesteer, wang2025truthflow, rodriguez2025controlling} or working on a more faithful manifold by rotating activations on the hypersphere \cite{you2026spherical, pham2024householder}. 
\noindent~Despite this progress, existing methods share two limitations. First, most of the methods and their implicit objective are heuristic rather than \emph{derived} from a well-posed optimization problem without formal justification for the objective itself. Second, these steering signals are \emph{query-independent}. A steering vector that helps in-distribution prompts can make out-of-distribution (OOD) inputs wrong or gibberish, and our experiments confirm that several strong baselines improve TruthfulQA~\cite{lin2021truthfulqa} at the cost of severe MMLU~\cite{hendrycks2020measuring} and GSM8K~\cite{cobbe2021gsm8k} degradation. We address both gaps with \Cobras, which casts steering as a Schr\"{o}dinger Bridge, yielding both a principled \emph{derivation} of the steering objective and an inherently query-adaptive steering direction, along with its OOD capability.

Overall, our contributions are: (i) we show that the Schr\"{o}dinger Bridge on the activation hypersphere provides a principled derivation of the implicitly widely used log-density-ratio steering objective; (ii) we propose Conditional Optimal Bridge for Riemannian Activation Steering (\Cobras), a practical algorithm that combines entropic optimal-transport potentials and an inherent gating mechanism to produce query-adaptive steering directions; and (iii) we demonstrate that \Cobras consistently outperforms prior baselines on helpfulness, truthfulness, and detoxification while avoiding the out-of-distribution degradation observed in existing methods.


\paragraph{Organization.}Section~\ref{sec:related} discusses the related work. Section~\ref{sec:bg} elaborates on the background, notation, setup, and the assumption of the manifold where LLM activations exist. Section~\ref{sec:method} then formulates the steering problem and its objective as the Schr\"{o}dinger Bridge instance, solves this problem, derives the objective for steering, and proposes \Cobras. Section~\ref{sec:exp} then does extensive experiments for both in-distribution and OOD and shows that \Cobras outperforms baselines in both settings. Finally, Section~\ref{sec:conclusion} concludes the paper, discussing limitations and future work. The proof of technical results and additional experiments are reported in the appendix. 


